\documentclass[11pt]{article}
\usepackage[utf8]{inputenc}
\usepackage[T1]{fontenc}
\usepackage{lmodern}
\usepackage[margin=1in]{geometry}
\usepackage{csquotes}
\usepackage[style=authoryear,backend=biber,sorting=nyt,maxbibnames=99,giveninits=true,uniquename=init]{biblatex}
\usepackage[hidelinks]{hyperref}
\addbibresource{references.bib}
\begin{document}
\section*{Bioreactor Design and Research Gaps}

\subsection{Overview}
The design of edge-of-field denitrifying bioreactors and their companion treatment structures has matured to the point where discrete performance demonstrations exist, yet the field is still defined more by open questions than by settled design rules. The literature synthesised here converges on several interlocking research frontiers: (i) how to size and configure bioreactors robustly, including for undersized systems and hydrologically flashy sites; (ii) how to select, amend, and manage reactive media---woodchips, corncobs, and specialised phosphorus-sorbing materials---to extend performance and lifespan; (iii) how to constrain the pollution-swapping and unintended by-products (greenhouse gases, dissolved organic carbon, phosphorus release, methylmercury) that accompany reductive treatment; (iv) how contaminants such as nitrate, phosphorus, pesticides, pathogens, and antibiotic resistance genes are transported through drainage inlets and macropore networks; and (v) how longer-term, multi-site monitoring and improved models can turn site-scale results into transferable design guidance. A recurring meta-message across studies is the call for standardised design procedures, broader performance-metric reporting, and long-term data collected across climates, landscapes, and maturation stages so that watershed-scale nutrient-reduction goals can be met \parencite{Christianson2021,Maxwell2022,Bell2015}. The subsections below organise these gaps by topic while retaining the specific hypotheses, methods, and conditions the primary studies identified as still needing resolution.

\subsection{Bioreactor sizing, scaling, and design procedures}
A central frontier is the reconciliation of site-level effectiveness with watershed-scale load removal. \textcite{Christianson2026} argue that effort should be directed toward better balancing bioreactors that are as effective as practical at the site level while still removing as much nitrogen load as possible, note that it would be valuable to assess the performance of undersized bioreactors in larger drainage areas, and suggest that greater flexibility in design processes may help scale the practice to those larger areas; they further caution that additional empirical data are needed to strengthen the loading-density relationship and that performance data from full-size bioreactors beyond Illinois are required to refine prediction models before they are used for design. In parallel, \textcite{Dougherty2020b} propose that particularly flashy drainage sites could consider sizing up bioreactor designs to meet performance goals, and repeatedly call for exploring simpler, standardised ``plug-and-play'' design approaches, positioning their flow analysis as a step toward a simplified design procedure. Standardisation is echoed at the practice-policy level by \textcite{Christianson2021}, who recommend refining the Conservation Practice Standard and streamlining the design process.

The adequacy of residence-time (T)-centric design is itself questioned. \textcite{deOliveira2023} advocate additional field-scale assessment of baffles in relatively wide bioreactors to increase hydraulic loading while maintaining hydraulic efficiency, call for more tracer testing at full-size bioreactors of varying designs coupled with nitrate monitoring and modelling, and specifically flag the weaknesses of current T-centric design approaches as a target for optimisation---including field-scale tracer testing of non-standard media. Foundational design-procedure validation remains incomplete: \textcite{Christianson2013} call for further field-scale validation of the drainage-bioreactor design procedure against performance data from bioreactors of different designs, and \textcite{Ghane2015} recommend validating hydraulic models at additional denitrification-bed sites and using them to develop optimal design guidelines, alongside using additional conservative tracers such as isotope tracers to better characterise bed hydraulics. Region-specific parameterisation is also outstanding; \textcite{Povilaitis2020} note that flow rate, seasonality, microbial community, and other limiting factors must be evaluated to obtain reliable bioreactor design parameters for Lithuania, and \textcite{Ghane2019} present bromide data intended to support the development of further sorption experiments underpinning tracer-based design.

\subsection{Reactive media, carbon sources, and substrate innovation}
Beyond conventional woodchips, several alternative fill media and carbon-management strategies remain incompletely characterised. \textcite{Law2023} encourage more field studies of corncobs as a bioreactor medium, including longer-term research capturing spatial and temporal variability for robust cost analysis, evaluation of the nitrogen removal rate range and whether a lower maximum hydraulic residence time (HRT) is needed, investigation of unintended pollution swapping such as nutrient and pesticide leaching, and assessment of greenhouse-gas production and biophysical properties under different environmental and operating conditions. \textcite{Porter2015} call for research into organic-carbon availability at different depths while accounting for moisture content, noting that if a link between moisture, carbon availability, and performance can be established, future designs could be altered---for example by using different substrates or by raising or lowering internal water levels---and suggest extending short-term microbial-diversity studies over longer periods to reveal recurring annual patterns in alpha diversity.

Carbon dosing and synthetic-carbon supplementation are active but unresolved. \textcite{Rivas2020} identify the specifics of adding readily available carbon---source, frequency, and amount---to optimise dissolved reactive phosphorus (DRP) and performance under highly dynamic inflow, while cautioning that more work is needed to minimise excess dissolved organic carbon (DOC) in outflow and potential biofilm clogging when carbon dosing is used. \textcite{Zhang2021b} identify continued study of operational parameters---especially HRT---the hybrid application of synthetic carbon sources, the design of low-temperature-tolerant reactors, and the mixing of synthetic carbon sources for complex water bodies as worthwhile directions. Microbial amendment strategies also require scale-up testing: \textcite{Jeglot2022} recommend that bioaugmentation, possibly combined with biostimulation, be tested under replicated pilot-scale conditions, with further testing to provide empirical and statistical evidence for upscaling to full-scale bioreactors treating agricultural drainage water. At the mechanistic end, \textcite{Ma2021} call for investigation of the inhibition and denitrifier-enhancement mechanisms caused by heated tourmaline and for field-scale tests of such bioremediation strategies for subsurface drainage water.

\subsection{Envelopes, filter aggregates, and structural materials}
A distinct materials frontier concerns drain envelopes and structural fill. \textcite{Kaboosi2012} recommend evaluating envelopes with different combination ratios of rice husk, sand, and gravel and conducting field experiments under real conditions before drawing conclusions on rice husk as a drain envelope, while \textcite{Navabian2017} stress long-term field experiments to determine the environmental and technical function of a combined rice-husk-and-sand envelope, given the degradation potential of rice husk and its tendency to create acidic conditions in the envelope medium. For aggregate sizing, \textcite{Byrne2023} recommend that aggregates ranging from 0.7 to 19~mm be tested in situ in clay-textured soils, with the identified optimal subset subsequently field-tested. Synthetic envelopes require material-by-material evaluation: \textcite{Byrne2023b} note that needle-punched non-woven geotextile envelopes need further testing, that thicker synthetic envelopes with performance comparable to aggregates may be required, that each synthetic envelope must be tested individually because of varying physical properties, and that filter-sock-plus-aggregate combinations should not be recommended until further research is carried out. Structural concrete alternatives also remain untested at scale; \textcite{Wang2016c} call for research into the practical application and durability of fly-ash- and coal-gangue-based concrete under field conditions with high groundwater levels.

\subsection{Phosphorus-removal media and filter design}
Phosphorus-sorption materials (PSMs) constitute one of the most heavily populated research-gap areas, spanning material selection, longevity, regeneration, and hydraulic integration. General selection criteria are laid out by \textcite{Li2020}, who argue that adsorbent choice must account for phosphorus concentration, interfering ions, temperature, flow rate, and sediment loads, prioritise selectivity for phosphorus at neutral pH, withstand seasonally fluctuating temperature, redox condition, flow rate, and sediment load without sacrificing removal efficiency, be testable for reusability and ease of recovery, and be economical for large field-scale application. \textcite{Mendes2022} add that the apparently synergistic retention mechanism and performance of filter media at high flow rates remain unknown, and that recycling and regeneration of spent filter media should be verified and pursued.

Iron- and steel-based media dominate the specific proposals. \textcite{Allred2017} call for optimising a modified iron-oxyhydroxide regeneration process using dilute hydrochloric acid followed by a NaOH flush, developing practices for collecting and treating the resulting acid/NaOH effluent, testing larger iron-oxyhydroxide systems capable of treating a greater percentage of total drainage flow, and evaluating low-cost goethite sources such as bog iron ore. \textcite{Allred2014} recommend more laboratory testing of zero-valent iron and related media to determine minimum contact time for phosphate removal, higher-flow-rate column tests, rigorous adsorption-isotherm tests to identify removal mechanisms, and field-scale filter treatment studies as a follow-on. \textcite{Hua2016} note that steel by-products may also remove methylmercury and may be regenerated by alkaline washing to recover adsorbed phosphate. \textcite{Sellner2019b} suggest connecting a steel by-product filter to a woodchip bioreactor only after the initial bioreactor flush (e.g. 1--2 weeks) to minimise DOC interference, and mixing fine steel particles with other media to increase hydraulic conductivity and reduce clogging. \textcite{Wang2020d} call for optimising the relationship among removal efficiency, retention time, steel-slag mass, and cost when designing a P-removal structure, while \textcite{Shedekar2020} plan to replace spent slag with a lower-cost gravel--metal-shavings mixture, note that future structures should use a thicker slag bed and reduced length to improve uniform flow distribution, and emphasise controlling sediment influx---especially from bank erosion---as a design consideration. \textcite{Penn2020} recommend PSMs that remove phosphorus via ligand exchange onto aluminium and iron oxide/hydroxide minerals (mine-drainage residuals, manufactured PSMs, aluminium-coated steel slag) and note that such materials suit tile-drainage treatment only where structures are designed for high-frequency media replacement---every 4--6 months in their study.

Fly-ash and pellet media raise saturation-timing questions. \textcite{Li2017} stress that saturation time must be quantified before field application so that fly-ash pellets can be replaced before they shift from phosphorus sink to source, and \textcite{Li2018b} direct follow-up research toward the hydraulic parameters of fly-ash pellets combined with recovery, and toward constraints such as poor saturation time, high outflow pH, and possible field desorption. Filter-media blends more broadly remain unoptimised: \textcite{Agrawal2011} call for determining the feasibility of filter installation and long-term field-scale performance, and \textcite{Agrawal2011b} report a need to determine the optimal media blend and configuration, to identify which constituent(s) offload phosphate and for how long before saturation, and to design larger or new filters capable of accommodating high drainage flows with sufficient contact time. Contact-time resolution itself is flagged by \textcite{Kumari2025}, who suggest exploring shorter contact-time intervals (e.g., 30~min) to capture the rapid initial adsorption phase. Anisotropic-medium behaviour warrants further study per \textcite{Kankarej2023}. Integrated woodchip-plus-PSM designs need multi-scale evaluation: \textcite{Husk2018} call for assessing mixed-media phosphorus-removal effectiveness under additional climate, soil, and cropping conditions over multiyear periods and at different scales to develop predictive design models, evaluating alternative or multiple media placements within bioreactors, examining effective P-removal lifespan and reuse/regeneration/recycling potential, and studying operational parameters such as pH, biochemical oxygen demand, and conductivity, including microbially mediated processes and removal of agrochemicals and microbial contaminants.

\subsection{Novel phosphorus dynamics and monitoring techniques}
Emerging mechanisms of phosphorus interaction with woodchips remain to be elucidated. \textcite{Zhang2024b} call for focused experiments on the processes leading to enhanced orthophosphate uptake by woodchips undergoing anoxic--(anaerobic--)oxic cycling, phenotypic identification of polyphosphate-accumulating organisms in woodchip biofilms to confirm enhanced-biological-phosphorus-removal-like mechanisms, and determination of whether anaerobic conditions are necessary to promote the observed uptake. Redox-driven phosphorus release is a recurring monitoring gap: \textcite{Bailon2025} call for evaluating bioreactor start-up flushes with immediate monitoring, aging bioreactors, and total phosphorus alongside higher-frequency DRP monitoring, and for further assessment of redox-driven P release and the use of drawdown plates. \textcite{Rivas2020} likewise call for continued multi-season monitoring to confirm the longer-term DRP removal potential of woodchip bioreactors and further study of the dynamics and factors determining DRP removal. Passive-sampling monitoring itself needs development: \textcite{Li2020b} call for a more reactive hybrid resin adaptable to a wide range of temperature and flow rate, investigation of particulate-phosphorus effects on passive-sampling estimates, mesh-bag materials with reduced biofouling for long-term deployment, and extension of the technique to large-scale watershed and river-basin monitoring. Fundamental characterisation of dissolved organic phosphorus is sought by \textcite{Liu2024b}, who call for investigating the roles of the origin and management of organic matter and organic P on P lability and DOP leaching, and the bioavailability of DOP in recipient waters.

\subsection{Pollution swapping and unintended by-products}
Reductive treatment introduces trade-offs that are only partly quantified. \textcite{Christianson2021} name reducing pollution swapping---especially at start-up---and improving understanding of the removal of additional pollutants, compounds, and biologicals as key research areas. Greenhouse-gas emissions are a prominent concern: \textcite{Brunton2024} recommend further greenhouse-gas research on bioreactors, fuller assessment of full-size bioreactor methane dynamics, examination of the impacts of bioreactor age and associated woodchip properties on GHG consumption and release, evaluation of ``hot moment'' GHG production, additional monitoring of methane production at new bioreactors, longer-term analyses paired with in situ woodchip physical and chemical property measurements, and further field-scale investigation of soil-cover benefits and trade-offs; they also suggest testing sophisticated automated control structures for real-time management based on outflow oxidation-reduction potential. \textcite{Strock2025} note that, because a 2-h HRT has been shown to increase greenhouse-gas emissions, additional research on this design consequence is needed, along with work on the longevity of P-removal materials, optimisation of performance under cold conditions, monitoring of DOC export, and strategies for managing P-treatment materials to minimise unintended phosphorus release. \textcite{Maxwell2022} suggest that minimising negative by-products such as mercury methylation could be based on an outflow nitrate-concentration design criterion rather than a mass metric. Nitrite and long-term P transport are flagged by \textcite{Bell2015}, who call for measuring nitrite concentrations in future bioreactor studies and for long-term evaluation of phosphorus transport from bioreactors. Carbon-export accounting more broadly is a monitoring gap: \textcite{Schilling2023} call for improved monitoring of aqueous total-carbon export for accurate carbon-sequestration accounting, multiple years of data to characterise carbon export under variable climate, and note that reducing water throughput via water-level control structures may minimise dissolved-inorganic-carbon loss.

\subsection{Longevity, clogging, subsidence, and construction}
The operational lifespan and physical stability of treatment structures are widely cited as under-studied. \textcite{Kohn2023} state that additional research is needed before widespread implementation, that bioreactor-function longevity is an important consideration, that estimating replacement interval and cost depends on key unknowns (performance, lifespan, and actual producer cost of agricultural residues), and that using a pre-measured feedstock volume per bioreactor with consistent placement could mitigate construction-method effects. Chemical clogging is mechanistically unresolved: \textcite{Guo2020} call for exploring the mechanism of chemical clogging and developing simulation models for technical prevention, a need echoed by \textcite{Guo2021}. Subsidence monitoring is a new methodological frontier; \textcite{Christianson2025b} suggest applying their LiDAR method at additional sites---especially following construction---to build subsidence time series, addressing whether subsidence is faster immediately after construction, collecting 3D surface data immediately post-construction, and assessing the practicality and benefit of targeted recharge near the inlet. Sediment management at construction is a design frontier for catchment-scale systems: \textcite{Feyereisen2023} plan to investigate the opposing effects of bypass flow (decreasing removal) and nitrate removal rate (increasing removal) in a sediment-exclusion system so managers can understand the net impact, recommend installing data-logging pressure transducers in each bed from project outset, and advise a pre-construction investigation of sedimentation sources with sediment-prevention measures included in design. Network-scale diagnostic methods also need extension; \textcite{Henine2021} recommend using an air compressor with higher injection pressure and airflow to extend a diagnostic method to much larger subsurface-drainage networks, monitoring electrical-resistivity variations with time-domain frequency electromagnetic methods and precise geolocation, and developing a safe air-injection system operable directly from the drain outlet.

\subsection{Pumped and energy-assisted bioreactor designs}
Pumping to increase hydraulic loading is an emerging but incompletely validated design avenue. \textcite{DíazGarcía2025} call for additional assessment of different pumping configurations and bioreactor designs to contextualise benefits and drawbacks, note that small but consistent post-batching temperature increases merit further study, observe that a power-storage unit could have increased flow treatment, identify operating pumps continuously for 24~h and doubling the pumping window to six months as hypothetical ways to increase hydraulic loading, and caution that pumping should perhaps be considered only in the most suitable applications. \textcite{Hartfiel2024} call for exploring solar-powered pumped bioreactors to reduce operating and unit costs, verifying pumped-bioreactor performance before policy implications can be realised, and optimising system performance while reducing unit cost. Pumped constructed-wetland-plus-woodchip-bioreactor systems similarly require further work: \textcite{Pugliese2020} call for testing that analyses the concentration and proportion of particulate and dissolved nutrient forms during highly variable, intense discharges, a cost-effectiveness analysis of the synergistic surface-flow-wetland-plus-woodchip-bioreactor combination versus standalone systems, and longer monitoring to define the pumping effect on total-phosphorus retention.

\subsection{Contaminant transport through drainage inlets and macropores}
Understanding how contaminants enter and move through drainage networks is explicitly under-studied. \textcite{Li2017b} recommend additional research on how nitrate, phosphorus, and pesticides are transported through drainage inlets, studies of inlet performance under different debris from different crop systems, and modification of the flow equations to represent actual inlet performance under debris conditions. The definition of a bioreactor's effective contributing area is itself uncertain; \textcite{Akay2008} question whether the contributing area should vary for water, tracers, and reactive contaminants. Sediment behaviour within tiles is likewise unresolved: \textcite{Stone1997} note that the relative importance of bacteria versus solution-chemistry changes on flocculation of tile sediment across a range of turbulent shear stresses remains to be determined.

Pesticide transport modelling forms a coherent cluster of gaps. \textcite{Wendell2024} propose implementing algorithms describing particle-mediated colloidal transport, delineating upland areas and floodplains together with pesticide accumulation when surface runoff infiltrates or evaporates, separately accounting for macropore flow when calculating pesticide leaching, and continuous improvements in modelling and observation to reduce uncertainty in pesticide degradation and transport predictions. \textcite{Zehe2001} argue that convection--dispersion or transfer-function process models must be improved to include pesticide transport by mobile particles and that application guidelines should include soil-specific interactions. \textcite{Fox2007b} identify the need for a quasi-two-dimensional procedure to distribute pesticide across the soil profile (rather than the current one-dimensional approach), site-specific calibration of the enrichment-factor parameter, more sophisticated accounting for macropore/subsurface-drain interconnectivity in chemical routing, and user awareness of model limitations. \textcite{Novak2003} call for better field-scale evaluation of metolachlor retention describing slow immobilisation, and for apparent-dissociation-kinetics processes to be incorporated in registration and management models, especially for persistent pesticides. \textcite{Dousset2004} argue that field-scale monitoring of herbicide losses in drains is necessary to compare molecule behaviour across soil types over at least two planting periods to reliably assess leaching risk, and \textcite{Leu2004} call for more quantitative functions relating available field characteristics to a field's vulnerability to lose herbicides. Analytical and marker methods also need development: \textcite{MUIR1978} note that persistence of triazine degradation products could be modelled under better-controlled conditions and that methods for determining hydroxy-triazines in water are needed for routine monitoring, while \textcite{Heaven2012} suggest 3,5-dichloro-2-methoxybenzoic acid could serve as a chemical marker of dicamba pollution and stress identifying contaminant sources. Continued monitoring is advised by \textcite{Gordon2005} given widespread linuron use in carrot and potato production. Over-winter conditions are singled out by \textcite{Siimes2006}, who call for detailed study of soil freezing/thawing and snowmelt effects on herbicide fate and note their dataset will test pesticide-fate models in over-winter conditions.

\subsection{Pathogens, microbial ecology, and antibiotic resistance}
Microbial contaminant transport and treatment span several open questions. \textcite{Soupir2018} call for investigation of pathogen removal under lower HRTs, examination of both \textit{Salmonella} and \textit{E.\ coli} removal at lower HRT values, exploration of the removal mechanism to inform longevity, and investigation of microbial removal at higher flows. The ecology of indicator bacteria in treatment wetlands is unresolved: \textcite{Stott2023} call for determining the relative survival and growth of faecal \textit{E.\ coli} strains versus non-faecal ``\textit{E.\ coli}-like'' \textit{Escherichia} (cryptic clades) to understand wetlands as sources of \textit{E.\ coli} and downstream health risks, and to confirm or reject whether wetland and aquatic habitats support extended persistence and growth. \textcite{Truhlar2015} ask whether curli expression responds to manure-application method and call for characterising downstream risks associated with curli-producing versus non-producing \textit{E.\ coli}, especially regarding foodborne illness. Long-term pathogen leaching following manure application needs study: \textcite{Malik2004} call for additional temporal and long-term studies on the survival and leaching of pathogens and indicators into subsurface water before firm recommendations can be made for swine-manure application. Fecal-source tracking methods require corroboration; \textcite{Gottschall2013} recommend using other markers---such as pharmaceuticals and personal-care products and fecal indicator bacteria---in tandem with sterols, and using multiple sterol ratios for cross-validation. Irrigation-related pathogen risks lead \textcite{Ding2021} to strongly recommend disinfection of wastewater effluents used for irrigation, lower winter irrigation rates to slow pathogen movement through wet, cold soil, and year-round supplemental disinfection. Broader catchment-scale microbial priorities are set out by \textcite{Collins2007}, including quantifying off-stream water, shade, and shelter sources for cattle; understanding cattle attraction to waterways to prioritise fencing; experimental studies across soils, slopes, and vegetation to improve understanding of microbial entrapment in riparian buffers; evaluating constructed wetlands, vegetated ditches, and farm ponds for treating microbes in stormwater and effluent; assessing the costs of advocated mitigation measures; better understanding water as a pathogen vector for human infection; tracing pathogen strains causing cattle re-infection; and examining links between cattle productivity and drinking-water microbial quality.

Antibiotic resistance genes (ARGs) are an explicitly emerging frontier. \textcite{Smith2019b} call for better understanding of the diversity of ARGs, their mechanisms of persistence and dissemination, and long-term environmental effects, and recommend targeted approaches such as qPCR to complement metagenomic characterisation. \textcite{Rieke2018} note that additional studies are needed to link the presence of resistance genes with host identification and transfer mechanisms.

\subsection{Companion practices: saturated buffers, wetlands, and controlled structures}
Frontier questions extend to practices deployed alongside or instead of bioreactors. For saturated buffers, \textcite{Ghane2025} recommend determining actual performance using a paired-field approach, further economic analysis of the buffer-flow component's cost-effectiveness, and determining average nitrate-load removal by the controlled-drainage and buffer-flow components under varying conditions; \textcite{Jaynes2019b} highlight the need to establish and monitor additional saturated buffers with multipipe configurations. Treatment-wetland performance limits are a persistent theme: \textcite{Tanner2005} call for longer-term studies of sustainable nutrient-removal rates across climatic conditions, process-based laboratory and mesocosm investigations to supplement field studies, determination of whether organic-carbon availability limits denitrification and whether organic-matter supplements or strategic biomass cutting increase removal, evaluation of different flow regimes using hydraulic residence-time distributions and dynamic water-mass-balance models, and further information on attenuation of other contaminants (including pathogens), net greenhouse-gas emissions, and biodiversity effects. \textcite{Tanner2003} recommend continued monitoring across seasons and years, intensive piezometer monitoring to quantify nutrient mass flows, studies of component processes to identify factors limiting nutrient removal, and supplementation of organic-carbon substrates with slow-release materials such as sawdust or straw. \textcite{Tanner2005b} note that further measurement and modelling are needed to estimate the frequency and composition of additional inflow sources, that performance might be optimised through different plant communities, soil types, organic-carbon amendments, or subsurface-flow zones, that longer-term studies are required to tease out event-to-event and year-to-year loading effects, and that larger wetland areas and/or enhanced areal denitrification rates would be needed to substantially reduce nitrogen export from pulsed drainage flows. Land-treatment wetland systems (FILTER) require subtropical adaptation per \textcite{Gardner2001}, who intend to increase applied-effluent nitrogen to representative values (about 25--30~mg~L\textsuperscript{-1} total nitrogen) and to clarify pathogen behaviour by tracking bacteriophage MS-2 after spiking. Controlled and flow-regulating structures also have open questions: \textcite{Popek2021} recommend further research on improving the connection of regulator elements to limit effective flow sizes, and \textcite{Schilling2023} note that water-level control structures may reduce dissolved-inorganic-carbon loss. Water and sediment control basins (WASCoBs) warrant study of pipe-riser type effects on nutrient and sediment transport, event-scale sediment and phosphorus trapping within ponded areas, berm-area-plus-pipe-riser interactions, and the impact of directly routing larger sub-basins to the watershed outlet \parencite{Gupta2019}. Nutrient-loss management and drainage-inlet flow control are further addressed by \textcite{Nédélec2010}, who note the need for additional work on soil-surface clogging, storage-area geometry, and internal-drain discharge capacity, experimental validation of flow, sediment, and pollution-mitigation efficiency, and study of detrimental effects such as bio- or soil accumulation of pollutants.

\subsection{Riparian, hyporheic, and effluent-management processes}
Subsurface attenuation processes adjacent to drainage require broader characterisation. \textcite{Puckett2005} call for more intensive study of hyporheic-zone processes at the study site, establishment of sufficiently wide buffer zones to guarantee adequate residence times, and further work on hydrogeologic controls on nitrate transport through riparian zones across a wider range of settings. Effluent- and manure-management gaps include \textcite{Monaghan2007}, whose ongoing work indicates low-rate effluent irrigation can reduce direct transfers of effluent phosphorus and bacteria through mole-pipe drainage by up to 95\%, and \textcite{Gordon2005}, who call for research on the microbiological safety of liquid-hog-manure practices and long-term (>10~yr) field trials with repeated manure applications to determine water-quality impacts relative to standard fertility practices. Nutrient-application optimisation is targeted by \textcite{Farahani2020}, who call for determining nitrogen and phosphorus application rates and timings to reduce concentrations, and by \textcite{Singh2002}, who recommend splitting a 60~kg~N~ha\textsuperscript{-1} basal dose into three 20~kg~N~ha\textsuperscript{-1} doses at 10-day intervals. Tillage and soil-biology effects on phosphorus leaching remain open: \textcite{Ulén2018} call for further monitoring to confirm any trend of reduced efficiency, long-term studies of untilled and shallow-tillage plots, investigation of biological factors such as earthworm populations, and long-term monitoring to quantify phosphorus-leaching effects, particularly for fields treated in several ways.

\subsection{Selenium and other trace-element legacies}
Trace-element contamination associated with agricultural drainage in the western United States remains a long-horizon research and monitoring commitment. \textcite{BENSON1993} outline continued biological and groundwater monitoring, continued research to identify long-term (10--20 year) strategies for selenium dissipation and removal from surface soils, annual site-management planning submitted to the regulatory board, a pilot-scale demonstration of microbial volatilisation in a former cattail area using a 0.8~ha test plot with irrigation and tillage, and evaluation of a no-action management alternative. \textcite{SAIKI1987} note that additional studies are planned or underway to determine whether fish and other aquatic animals elsewhere in the San Joaquin Valley are also contaminated or adversely affected by selenium. Ecotoxicological monitoring of other agricultural contaminants is likewise recommended: \textcite{Campbell2023} call for investigating the mechanisms behind decreased cerebellum size, a more comprehensive analysis of predation, competition, and other ecological factors at field-collection sites, further examination of the relationship between imidacloprid and imidacloprid-olefin brain levels under field conditions, and regular monitoring of neonicotinoid levels in waterways near agricultural activity with tighter environmental restrictions to protect vulnerable non-target organisms such as amphibians. \textcite{Dreyfuss2018} note that findings for \textit{Galba truncatula} require reviewing the biogeographical status of the species and considering conservation measures given its role as intermediate host in the \textit{Fasciola hepatica} cycle.

\subsection{Cross-cutting priorities: long-term monitoring, metrics, and models}
Synthesising across topics, several priorities recur regardless of practice. First, long-term and multi-site data: \textcite{Bell2015} call for more long-term studies across a variety of landscapes and more field-scale performance data from bioreactors of different designs and multiple Midwest locations; \textcite{Christianson2021} call for more long-term studies to advance bioreactor technologies, and \textcite{Bailon2025} recommend additional monitoring of a variety of bioreactor designs under different climates as treatment systems mature. Second, standardised and expanded performance-metric reporting: \textcite{Maxwell2022} argue that broadly reporting a variety of performance metrics will help scale both achievable treatment and implementation in non-point-source applications. Third, temperature-sensitive and mechanistic models: \textcite{Christianson2021} note that nitrate-removal rate analysis across regional or seasonal temperature gradients would provide added insight and that recent bioreactor models should incorporate temperature sensitivity. Together, these cross-cutting needs frame the overarching research agenda---moving from site-scale demonstration toward validated, transferable design guidance capable of meeting watershed-scale water-quality goals across the diverse climates, soils, and hydrologic regimes in which subsurface drainage operates.

\printbibliography
\end{document}
